\section{Transcriptome Transformer with PPI-guided volumetric attention}
\label{sec:method_txt_vma}

The predictive model developed in this thesis combines the transcriptome-centred representation of the Transcriptome Transformer (TxT) with a sparse adaptation of the Volumetric Multimodal cross-Attention (VMA) introduced in GRAMformer \cite{koo2025transcriptome_transformer,cicchetti2026gramformer}. The two components address complementary aspects of the problem. TxT treats genes as an unordered collection of tokens and uses the expression profile of each subject to make self-attention patient-specific. VMA introduces a geometric term that measures the joint configuration of more than two representations. In the proposed integration, this term is not applied to separate input modalities. It is adapted to triples formed by a gene and one of its neighbours in a protein--protein interaction (PPI) network. The resulting architecture, referred to below as PPI-guided TxT-VMA, retains the dense TxT encoder and augments each self-attention layer with a sparse residual branch defined only over biologically supported edges.

\subsection{Input representation and gene embeddings}

For subject \(p\), the model receives an expression vector \(\mathbf{x}^{(p)}\in\mathbb{R}^{N}\), where \(N\) is the number of genes retained after feature selection. In the experiments reported in this thesis, the \(2{,}000\) genes with the highest variance are selected using the training partition only. A min--max transformation is then estimated independently for each selected gene on the same training partition and applied unchanged to validation and test subjects. Gene selection and scaling are repeated within every experimental split, preventing information from held-out subjects from affecting the input representation.

Each retained gene \(i\) is associated with a trainable embedding \(\mathbf{g}^{(0)}_i\in\mathbb{R}^{d_e}\). The original TxT formulation initializes these representations using node2vec embeddings learned from a PPI network \cite{koo2025transcriptome_transformer}. In the implementation used here, the embeddings are instead initialized from a zero-centred Gaussian distribution with standard deviation \(0.02\) and optimized jointly with the classifier. The biological network is introduced explicitly by the volumetric branch described below. This separation allows genes without a compatible pretrained embedding to remain in the transcriptomic input while making the role of the PPI prior directly auditable.

The genes are not assigned a sequential order. Their identity is represented by the learned embedding, whereas the expression value controls patient-specific attention. The input should therefore be interpreted as a set of gene tokens rather than as a sequence analogous to natural language.

\subsection{Dense TxT self-attention}

To keep the derivation readable, the patient, encoder-layer, and attention-head indices are omitted below whenever they are unambiguous. For one subject, layer, and head, let \(\mathbf{H}\in\mathbb{R}^{N\times d_{\mathrm{model}}}\) collect the current gene representations. The usual query, key, and value matrices are
\begin{equation}
    \mathbf{Q}=\mathbf{H}\mathbf{W}_Q,
    \qquad
    \mathbf{K}=\mathbf{H}\mathbf{W}_K,
    \qquad
    \mathbf{V}=\mathbf{H}\mathbf{W}_V,
\end{equation}
with head dimension \(d_h=d_{\mathrm{model}}/H\). If only these matrices were used, all subjects would enter the first encoder layer with the same gene embeddings. TxT makes the attention patient-specific through a transcriptome-based term inspired by untied positional encoding. The expression vector is projected separately for the query and key roles, producing \(\mathbf{Q}^{x}\) and \(\mathbf{K}^{x}\), and the resulting TUPE matrix is
\begin{equation}
    \mathbf{T}(\mathbf{x})
    =\frac{\mathbf{Q}^{x}(\mathbf{K}^{x})^{\top}}{\sqrt{2d_h}}.
\end{equation}
It does not encode the position of a gene in an arbitrary list. Instead, it modifies every gene--gene relation according to the expression values observed in the current subject. The complete dense attention is
\begin{equation}
\label{eq:txt_dense_logit}
    \mathbf{A}^{D}
    =\operatorname{softmax}_{\mathrm{row}}
    \left(
        \frac{\mathbf{Q}\mathbf{K}^{\top}}{\sqrt{2d_h}}
        +\mathbf{T}(\mathbf{x})
    \right),
    \qquad
    \mathbf{O}^{D}=\mathbf{A}^{D}\mathbf{V}.
\end{equation}
Attention is therefore dense: every retained gene can exchange information with every other retained gene. The outputs of the heads are concatenated, projected to the embedding dimension, and processed by a position-wise feed-forward network. Residual connections and layer normalization are applied after both the attention and feed-forward sublayers.

\subsection{From GRAMformer to PPI-guided volumetric attention}

GRAMformer replaces pairwise multimodal cross-attention with a score derived from the volume spanned jointly by a query and multiple modality-specific keys \cite{cicchetti2026gramformer}. Given a matrix \(\mathbf{Z}\) whose columns contain these vectors, the corresponding volume is obtained from the determinant of the Gram matrix \(\mathbf{Z}^{\top}\mathbf{Z}\). A small volume indicates that the vectors lie close to a common lower-dimensional subspace, whereas a larger volume indicates greater joint linear independence.

The setting considered in this thesis is not conventionally multimodal: every token corresponds to a gene from the same expression profile. The volumetric principle is therefore adapted to the local topology of a PPI graph \(\mathcal{G}^{(r)}=(\mathcal{V}^{(r)},\mathcal{E}^{(r)})\), induced at each experimental repetition \(r\) by the genes selected from the training data. The edge set is obtained from the high-confidence HIPPIE network by retaining interactions with confidence score at least \(0.73\) and with both endpoints in the selected gene set. Each undirected interaction is expanded into two directed edges. The confidence value determines whether an interaction is retained but is not used as an additional learned edge weight. Consequently, the volumetric branch depends on the topology of the induced PPI graph, while the dense TxT branch continues to operate over all gene pairs.

For an edge \(i\leftarrow j\), the local higher-order interaction combines the VMA query and anchor key of gene \(i\) with the key of its PPI neighbour \(j\). The score therefore depends jointly on the state of the destination gene and on the representation supplied by a biologically connected source gene.

\subsection{Patient-specific expression injection in VMA}

The volumetric branch receives expression explicitly, rather than relying only on the TUPE term. For each subject, expression is standardized across the valid selected genes:
\begin{equation}
    \widetilde{x}_i
    =
    \begin{cases}
        (x_i-\mu_x)/\sigma_x, & \text{if }x_i\text{ is valid},\\
        0,                    & \text{otherwise},
    \end{cases}
\end{equation}
where \(\mu_x\) and \(\sigma_x\) are the subject-specific mean and stabilized standard deviation computed only from finite measurements. No quantity in this operation depends on the other subjects present in the minibatch.

Three learned projections inject the standardized expression directly into the query, key, and value spaces:
\begin{equation}
\begin{aligned}
    \mathbf{Q}^{V}&=\operatorname{sg}(\mathbf{Q})+\widetilde{\mathbf{x}}\mathbf{p}_Q^{\top},\\
    \mathbf{K}^{V}&=\operatorname{sg}(\mathbf{K})+\widetilde{\mathbf{x}}\mathbf{p}_K^{\top},\\
    \mathbf{V}^{V}&=\operatorname{sg}(\mathbf{V})+\widetilde{\mathbf{x}}\mathbf{p}_V^{\top},
\end{aligned}
\end{equation}
where \(\operatorname{sg}(\cdot)\) denotes stop-gradient. Detaching the backbone tensors prevents the auxiliary branch from changing the shared query, key, and value projections through its own gradient path. The dense TxT branch remains fully trainable, while the VMA-specific projections, gates, and residual fusion are optimized from the task losses. This separation was adopted to stabilize the comparison with the unaugmented TxT backbone.

\subsection{L2-normalized volumetric score}

Raw Gram determinants are sensitive to the norms of their input vectors. To make the geometric term dimensionless, let \(\widehat{\mathbf{z}}=\mathbf{z}/\max(\lVert\mathbf{z}\rVert_2,\varepsilon)\). For edge \(i\leftarrow j\), the normalized local triple and its volume are
\begin{equation}
    \mathbf{Z}_{ij}
    =\left[
       \widehat{\mathbf{q}}^{V}_i,
       \widehat{\mathbf{k}}^{V}_i,
       \widehat{\mathbf{k}}^{V}_j
      \right]
    \in\mathbb{R}^{d_h\times 3},
\end{equation}
\begin{equation}
\label{eq:ppi_volume}
    \nu_{ij}
    =\sqrt{
       \max\!\left\{
       \det\!\left(\mathbf{Z}_{ij}^{\top}\mathbf{Z}_{ij}\right),0
       \right\}+\varepsilon
      }.
\end{equation}
The determinant is evaluated in single precision using the analytic expansion of the \(3\times3\) Gram determinant. Negative values caused only by floating-point error are clamped to zero before the square root.

The sparse attention logit is then
\begin{equation}
\label{eq:ppi_vma_logit}
    b_{ij}
    =
    \frac{(\mathbf{q}^{V}_i)^{\top}\mathbf{k}^{V}_j}{\sqrt{d_h}}
    -\beta\nu_{ij}
    +t_{ij}(\mathbf{x}),
    \qquad (i,j)\in\mathcal{E}^{(r)},
\end{equation}
where \(t_{ij}(\mathbf{x})\) is the corresponding entry of the TUPE matrix in Eq.~\eqref{eq:txt_dense_logit}, restricted to the PPI edge, and \(\beta\) controls the volumetric contribution. The negative sign assigns higher compatibility to triples with a smaller joint volume, following the alignment interpretation used by GRAMformer. Because \(\nu_{ij}\) is computed from unit vectors, it is not divided by \(\sqrt{d_h}\); standard scaling is retained only for the dot product.

The softmax is computed separately over the incoming neighbourhood of every destination gene:
\begin{equation}
    \pi_{ij}
    =\operatorname{softmax}_{j\in\mathcal{N}(i)}(b_{ij}).
\end{equation}
This segmented normalization allocates memory proportional to the number of retained PPI edges rather than to \(N^2\). It reduces the additional cost of the VMA branch, although the complete model still contains the original dense TxT attention and therefore retains its quadratic component.

\subsection{Neighbour--self message and residual fusion}

VMA produces a patient-specific graph contrast rather than a second dense attention output. For destination gene \(i\), define the neighbour aggregate \(\mathbf{m}^{N}_i\), the neighbour and self gates \(\mathbf{g}^{N}_i\) and \(\mathbf{g}^{S}_i\), and the resulting contrast \(\mathbf{r}_i\) as
\begin{equation}
\begin{aligned}
    \mathbf{m}^{N}_i
    &=\sum_{j\in\mathcal{N}(i)}\pi_{ij}\mathbf{v}^{V}_j,\\
    \mathbf{g}^{N}_i&=\sigma(\mathbf{W}_n\mathbf{q}^{V}_i),
    \qquad
    \mathbf{g}^{S}_i=\sigma(\mathbf{W}_s\mathbf{q}^{V}_i),\\
    \mathbf{r}_i
    &=\mathbf{g}^{N}_i\odot\mathbf{m}^{N}_i
      -\mathbf{g}^{S}_i\odot\mathbf{v}^{V}_i.
\end{aligned}
\end{equation}
Here, \(\sigma\) is the logistic sigmoid and \(\odot\) denotes element-wise multiplication. The difference represents the expression-aware information supplied by the PPI neighbourhood relative to the current state of the destination gene. Genes without an active retained neighbour receive a zero VMA message.

Before fusion, the VMA tensor is matched to the scale of the dense attention output independently for each subject. Let \(\operatorname{RMS}_p\) denote the root mean square computed over all heads, genes, and channels of one subject. The normalized graph contrast is
\begin{equation}
    \widetilde{\mathbf{R}}
    =
    \operatorname{sg}\!\left(
    \frac{\operatorname{RMS}_p(\mathbf{O}^{D})}
    {\operatorname{RMS}_p(\mathbf{R})}
    \right)\mathbf{R}.
\end{equation}
The detached scale factor avoids introducing a separate optimization path through the normalization statistic. More importantly, subject-wise RMS matching makes the prediction invariant to the composition of the minibatch.

Finally, each head has a learned bounded residual gate:
\begin{equation}
\label{eq:vma_fusion}
    \mathbf{O}^{(h)}
    =
    \mathbf{O}^{D,(h)}
    +\tanh\!\left(\gamma_h\right)
    \widetilde{\mathbf{R}}^{(h)}.
\end{equation}
All \(\gamma_h\) parameters are initialized to zero. The augmented model therefore begins from exactly the same forward mapping as its TxT counterpart, while training determines independently for each head the magnitude and sign of the volumetric correction. No dropout is applied to the sparse VMA weights in the main configuration.

This formulation differs from the original GRAMformer in three respects. First, the vectors in Eq.~\eqref{eq:ppi_volume} are not separate modalities but an anchor--neighbour configuration from a single transcriptomic representation. Second, VMA augments rather than replaces dense attention. Third, its value aggregation uses a neighbour--self contrast and a zero-initialized per-head residual gate, rather than averaging gated values from modality-specific encoders. GRAMformer provides the geometric principle; the PPI graph and TxT determine how that principle is instantiated for transcriptomic classification.

\subsection{Shared multitask prediction}

The encoder is shared across the three pairwise diagnostic tasks. After the final encoder layer, average pooling over the genes produces one subject representation
\begin{equation}
    \mathbf{z}^{(p)}=\frac{1}{N}\sum_{i=1}^{N}\mathbf{h}^{(p,L)}_i.
\end{equation}
Three task-specific multilayer perceptrons map this representation to the logits for AD versus MCI, AD versus CTL, and MCI versus CTL. Each head contains hidden layers of dimensions \(128\) and \(64\), batch normalization, LeakyReLU activations, and dropout. A task mask ensures that a subject contributes only to comparisons containing its observed clinical group. Thus, an AD subject contributes to AD versus MCI and AD versus CTL, while it is excluded from MCI versus CTL.

Let \(\mathcal{L}_{\mathrm{AM}}\), \(\mathcal{L}_{\mathrm{AC}}\), and \(\mathcal{L}_{\mathrm{MC}}\) denote the cross-entropy losses computed on the valid subjects for the three tasks. The joint objective is
\begin{equation}
\label{eq:txt_multitask_loss}
    \mathcal{L}
    =0.50\mathcal{L}_{\mathrm{AM}}
    +0.25\mathcal{L}_{\mathrm{AC}}
    +0.25\mathcal{L}_{\mathrm{MC}}.
\end{equation}
The larger weight assigned to AD versus MCI reflects its role as the primary and most difficult clinical comparison. Training batches are balanced across the three biological classes, without constructing synthetic expression profiles.

\subsection{Model configuration and checkpoint selection}

The configuration used for the main experiments contains one shared encoder layer, two attention heads, \(d_{\mathrm{model}}=d_e=64\), and a feed-forward dimension of \(256\). Encoder and head dropout are set to \(0.2\). The volumetric branch uses the L2-normalized score in Eq.~\eqref{eq:ppi_volume}, \(\beta=1\), subject-wise RMS normalization, per-head gates, detached backbone tensors, and zero VMA dropout. The encoder, embeddings, and prediction heads are optimized with learning rate \(10^{-4}\), while VMA-specific parameters use \(5\times10^{-4}\). AdamW is used with weight decay \(10^{-4}\). Gradient norms are clipped separately for the TxT and volumetric parameter groups at \(1.0\).

Training lasts at most \(200\) epochs. A numerical loss-stop terminates a run when validation loss remains above \(1.0\) for five consecutive epochs. Checkpoints are ranked using
\begin{equation}
    S_{\mathrm{val}}
    =0.70\operatorname{AUC}^{\mathrm{AM}}_{\mathrm{val}}
    +0.15\operatorname{AUC}^{\mathrm{AC}}_{\mathrm{val}}
    +0.15\operatorname{AUC}^{\mathrm{MC}}_{\mathrm{val}}
    -0.25\mathcal{L}_{\mathrm{val}}.
\end{equation}
The final predictor is the arithmetic mean of the class probabilities produced by the three highest-scoring checkpoints, constrained to be separated by at least three epochs. The test partition is not evaluated during training or checkpoint selection and is opened only once for the frozen ensemble.

Several ablations distinguish the contributions of the merged components. The TxT baseline removes the complete PPI-VMA branch. Setting \(\beta=0\) retains sparse expression-aware message passing but removes only the Gram-volume term from Eq.~\eqref{eq:ppi_vma_logit}. Setting the residual gates to zero at inference removes the complete VMA contribution without changing the trained dense path. Finally, disabling TUPE replaces the second term of Eq.~\eqref{eq:txt_dense_logit}, and the corresponding edge term in Eq.~\eqref{eq:ppi_vma_logit}, with zero. This last experiment tests whether direct expression injection into the VMA query, key, and value vectors can replace the patient-specific mechanism of the original TxT encoder.
